\section{Related Work}
\label{sec:related_work}
% Classical chess engines, alpha-beta, iterative deepening, time control

The problem of efficient search in chess engines has been studied extensively
over several decades~\cite{marsland1986review}. Early approaches were based on the Minimax algorithm,
which provides an optimal decision strategy under the assumption of perfect
play~\cite{yannakakis2018artificial}. However, due to the exponential growth of the game tree, Minimax quickly becomes infeasible even at moderate depths.

Alpha-Beta pruning was introduced as a major improvement over plain Minimax,
allowing large parts of the search tree to be pruned without affecting the
correctness of the result~\cite{knuth1975alphabeta}. Subsequent work showed that
the effectiveness of Alpha-Beta pruning depends crucially on the order in which
moves are explored. With perfect move ordering, the effective branching factor can be reduced dramatically, enabling much deeper searches within the same computational budget~\cite{marsland1986review}.

Iterative deepening~\cite{korf1985iterative} has emerged as a standard technique to improve move ordering
and to support time-controlled search. By repeatedly searching the game tree with
increasing depth limits, iterative deepening ensures that a valid best move is
always available while simultaneously refining the move ordering for deeper
iterations. Despite the apparent redundancy caused by re-searching parts of the
tree, it has been shown that the overhead is often compensated by improved
pruning efficiency~\cite{marsland1986review}.

A detailed analysis of Alpha-Beta pruning efficiency in chess engines is
presented by Nair~\cite{nair_improving_ab}. The paper investigates how different
move ordering heuristics and search enhancements affect the pruning behavior of
Alpha-Beta search. In particular, it highlights the strong interaction between
iterative deepening, transposition tables, and heuristic move ordering. The
results indicate that improved move ordering can significantly reduce the number
of explored nodes, even when additional search iterations are introduced.

Building directly on the observation that Alpha-Beta search is most effective
with strong move ordering, several refinements of the basic algorithm have been
proposed. Pearl introduced the Scout algorithm, which verifies promising moves with
minimal-window tests before evaluating them exactly~\cite{pearl1980scout}.
Marsland and Campbell transferred this idea to the Alpha-Beta framework in the
form of Principal Variation Search~\cite{marsland1982parallel}, and Reinefeld
later refined it into NegaScout, analyzing the conditions under which the
additional re-searches pay off~\cite{reinefeld1983negascout}. PVS is
particularly effective in combination with iterative deepening and
transposition tables, since both increase the probability that the first move
searched is indeed the best one~\cite{marsland1986review}.

Transposition tables have further improved search efficiency by allowing engines
to reuse previously computed results for positions reached via different move
orders~\cite{zobrist1970hashing}. By storing evaluations together with depth
information and bound types, transposition tables enable both direct cutoffs and
enhanced move ordering. Their effectiveness is especially pronounced in
combination with iterative deepening, where information from earlier iterations
can be reused almost immediately~\cite{warnock1988search}. Beyond single search
invocations, a persistent transposition table also allows information to be
carried over from one move to the next, effectively resuming previously
explored search trees instead of recomputing them from scratch.

Quiescence search addresses another well-known limitation of depth-limited
search, commonly referred to as the horizon effect. By extending the search in
tactically unstable positions, quiescence search produces more reliable
evaluations and has become a standard component of classical chess
engines~\cite{marsland1986review}.

Beyond exact pruning, engines also employ heuristic forward pruning, which
discards subtrees without proving them irrelevant. The most widely used variant
is null-move pruning: if passing a turn still leaves the position good enough to
cause a cutoff, the subtree is assumed to be safe to skip. The null move as a
search device is considerably older than this application of it, but the
recursive form with depth reduction that engines use today was established
experimentally by Goetsch and Campbell~\cite{goetsch1990experiments}. Unlike
Alpha-Beta this is not value-preserving, so it trades a small risk of overlooking
a line for a substantially smaller tree.

Closest to the question of this thesis is the work on how additional search
actually pays off. Junghanns et al.~\cite{junghanns1997diminishing} studied
diminishing returns in self-play and found the winning percentage per additional
ply to decrease with depth. They also report that this decrease is not visible
throughout: the high error rate of chess programs hides the effect for search
depths through nine plies, which is precisely the range in which the engine
studied here operates. Ferreira~\cite{ferreira2013impact} approached the
same relationship from the other direction and related search depth to Elo,
reporting gains per ply that shrink as depth grows.

These studies establish that the returns of additional search diminish, but they
measure it through game outcomes between engines of differing depth. What they do
not address is where the returns diminish \emph{within a single engine under a
clock}, at which depth its decisions stop changing, and which positions cause the
remaining error. This thesis takes that perspective: instead of relating depth to
playing strength, it relates the time budget to the quality of the individual
decision on a fixed set of positions, and separates the positions in which
additional search still changes the decision from those in which it does not.

While modern state-of-the-art engines incorporate additional techniques such as
bitboards, evaluation learning, and neural networks~\cite{silver2018alphazero},
the classical search principles discussed above remain fundamental. This work
builds upon these established ideas.